% Generated draft: requires mathematical and source review.
\documentclass[11pt]{article}
\usepackage[margin=25mm]{geometry}
\usepackage{fontspec}
\setmainfont{Latin Modern Roman}
\usepackage{amsmath,amssymb,mathtools,unicode-math}
\setmathfont{Latin Modern Math}
\usepackage{enumitem,needspace,xurl,hyperref}
\setlength{\parindent}{0pt}
\setlength{\parskip}{5pt}
\setlength{\emergencystretch}{3em}
\newcommand{\N}{\mathbb{N}}
\newcommand{\Z}{\mathbb{Z}}
\newcommand{\Q}{\mathbb{Q}}
\newcommand{\R}{\mathbb{R}}
\newcommand{\C}{\mathbb{C}}
\newcommand{\F}{\mathbb{F}}
\newcommand{\A}{\mathbb{A}}
\newcommand{\PP}{\mathbb{P}}
\newcommand{\E}{\mathbb{E}}
\newcommand{\eps}{\varepsilon}
\newcommand{\dd}{\,\mathrm{d}}
\begin{document}
\section*{UnsolvedMath ID: 11000215}
\textbf{Source problem code:} catalogue entry
\subsection{Definitions and mathematical statement}
A square-tiled translation surface \((X,\omega)\) is conventionally obtained from finitely many unit Euclidean squares by translation identifications of parallel edges, with \(\omega\) the induced holomorphic \(1\)-form. Equivalently it gives the corresponding branched cover of the standard flat torus, with the permitted branch data fixed by convention. A stratum \(\mathcal H(k_1,\ldots,k_r)\) records zero orders of \(\omega\); in the usual unmarked genus-\(g\) setting, \(\sum_i k_i=2g-2\).

The \(\mathrm{SL}_2(\R)\) action postcomposes the translation charts with a determinant-one linear map. Problem 8 asks to classify
\[
 \mathrm{SL}_2(\R)\cdot(X,\omega)
\]
for square-tiled surfaces in every stratum, and to describe their Teichmüller discs. The full orbit usually contains surfaces not tiled by unit squares, so the task concerns orbits \emph{containing} the square-tiled representatives, not a claimed action of all \(\mathrm{SL}_2(\R)\) on the discrete square-tiled set.

In the usual marked model a Teichmüller disc is generated by the affine deformations of one translation surface, with rotation factored out; its projection to moduli and its affine stabilizer encode further geometry. The source does not specify whether the requested description concerns the marked disc, its moduli image, stabilizer/Veech group, or all of these. It does not restrict to primitive torus covers or fix the number of squares. Area normalization matters because \(\mathrm{SL}_2(\R)\) preserves area, so arbitrary rescaling must not be silently included.

The source highlights \(\mathcal H(1,1)\), the genus-two stratum with two simple zeros, and says the problem was solved only for \(\mathcal H(2)\), the genus-two double-zero stratum, citing [HuLe1] and [McM4]. That ``only'' is a source-era assessment, not a new current-literature claim. The supplied summary reports continued orbit-classification study and an open general problem without further precise results.
\subsection{Short English statement}
Classify the determinant-one linear orbits of all square-tiled surfaces in each stratum and describe their Teichmüller discs, especially in \(\mathcal H(1,1)\). The source's \(\mathcal H(2)\) solved-case report is retained historically, with area, covering and disc-description conventions still to be fixed.
\subsection{Sources}
\begin{itemize}
\item Local UnsolvedMath registry, permanent ID \(11000215\), exact source code \texttt{catalogue entry}, Problem 8, ``Orbits of square-tiled surfaces,'' PDF page 259. No individual proposer or proposed year is supplied.
\item Farb, edited volume \emph{Problems on Mapping Class Groups and Related Topics}: \url{https://www.math.uchicago.edu/~farb/papers/mcgbook.pdf}. The volume has 329 explicitly typeset numbered items across chapters, repeated local numbering distinguished by PDF page.
\item Internal references [HuLe1] and [McM4] for the reported \(\mathcal H(2)\) case; full bibliography entries are absent and no author/title is guessed.
\item Extraction pdf-font-structure, accessed 2026-07-29. Source evidence NEEDS\_REVIEW presents an open problem and calls for a current primary-literature audit. Registry status open and classification OPEN-TRIAGE; supplied summary mentions \(\mathrm{SL}_2(\Z)\)-orbit/Teichmüller study and says the general classification is open, without exact scope or additional citations.
\item Dataset attribution: Ulam AI / UnsolvedMath Contributors, CC BY 4.0. Source rights NEEDS\_REVIEW: publicly accessible author/publisher copy, redistribution terms requiring review.
\end{itemize}
\end{document}
